\titlespacing*{\section}{0pt}{0.8em}{0.45em} \section{Final Summary}
VietFlood was designed as a working prototype to translate ad-hoc flood observations into structured reports that can be searched, evaluated and updated. It’s about reporting and coordination help, not about flood prediction. A citizen can file a report from a phone with location, description and optional photo evidence and that report is available for follow up instead of disappearing in chat history.

The supplied system consists of a NestJS backend divided into an API gateway, an authentication service and a reporting service, an Expo React Native mobile application for the citizens and relief flows and a Next.js dashboard for the administrators \cite{nestjsDocs,expoDocs,nextjsDocs}. Reports are persisted in PostgreSQL using TypeORM, evidence is uploaded to Cloudinary and saved as metadata, and report lookups can be sped-up with Redis caching \cite{typeormDocs,cloudinaryUploadDocs,redisDocs}. The gateway performs JWT based role checks for the segregation of citizen reporting, relief browsing and administrator review \cite{jones2015jwt}. The RabbitMQ message patterns are used for the service-to-service work, so the HTTP concerns remain in the gateway and the business logic in internal services \cite{rabbitmqConfirmsDocs,nestjsMicroservicesDocs}. This separation follows conventional microservice design recommendations, but is still small enough to be deployed as a prototype \cite{fowlerLewis2014microservices,dragoni2017microservices,thones2015microservices}

The prototype shows the complete cycle of reporting: report creation, archiving with evidentiary meta data, listing and retrieval using secure endpoints and change of status when approved by an administrator. In this version the relief role is deliberately read-focused. That decision maintains status updates in administrator control, since status is the primary trust signal that is shown back to other users.

In Chapter~5 the evaluation confirms the functional coverage by means of scenario walkthroughs and examination of the offered services and clients. VietFlood needs to be validated under real flood conditions with inadequate connectivity and more traffic before it can be considered for operational use.

\section{Recommended Next Steps}

The following enhancements should be based on genuine usage input on the reporting flow. If users are unable to quickly submit the minimum information (location, short description, and one photo), the system will not be useful in stressed settings.

From this prototype several concrete engineering tasks arise:
\begin{itemize}[leftmargin=1.2em]
    \item \textbf{Offline-first reporting:} Local drafts, queued uploads, retry logic and clear client states for weak networks. Control abuse and storage cost by compressing evidence on-device and enforcing server-side restrictions.
    \item \textbf{Data quality and de-duplication:} make DTO naming and validation more stable and add duplicate identification based on time, distance and visual resemblance to reduce review workload.
    \item \textbf{Improved review workflow:} add clustering views, increase filters, support lightweight triage notes, so that administrators can evaluate bursts of incoming reports.
    \item \textbf{Tracking hardening:} bind Socket.IO channels to authenticated identities and enforce permission rules in line with the planned operational policy.
    \item \textbf{Operations and security:} add health checks, monitoring and failure alarms, test refresh token rotation under concurrency, and improve deployment methods for repeatability and recovery \cite{humble2010continuousDelivery,kim2016devopsHandbook,dockerComposeDocs}.
\end{itemize}

In a longer project schedule VietFlood can also be expanded beyond reporting: integration with governmental warning channels, evidence verification collaborations and map-oriented aggregation for public information. The handling of those extensions should be done with care, as public facing status and location data brings privacy and safety problems that are out of the scope of this thesis.
